\DocumentMetadata{
  pdfversion=2.0,
  pdfstandard=ua-2,
  lang=en-US,
  tagging=on,
  tagging-setup={math/setup=mathml-SE,tabsorder=structure}
}
\documentclass[11pt,letterpaper]{article}
\usepackage[margin=0.68in,includefoot]{geometry}
\usepackage{newcomputermodern}
\usepackage{amsmath,amscd,mathtools}
\usepackage{tikz,adjustbox}
\providecommand{\square}{\mdlgwhtsquare}
\usepackage[hidelinks]{hyperref}
\hypersetup{
  pdftitle={Algebra First-Year Practice Exam A, Part 1},
  pdfauthor={University of Florida Department of Mathematics},
  pdfsubject={Graduate mathematics examination},
  pdfdisplaydoctitle=true,
  bookmarksopen=true
}
\setcounter{secnumdepth}{0}
\setlength{\parindent}{0pt}
\setlength{\parskip}{0.28em}
\setlength{\emergencystretch}{3em}
\setlength{\leftmargini}{2.15em}
\setlength{\leftmarginii}{2.65em}
\setlength{\leftmarginiii}{3em}
\pagestyle{empty}
\raggedbottom
\allowdisplaybreaks
\NewTaggingSocketPlug{block/list/label}{examproblem}{%
  \tagstructbegin{tag=Lbl}\tagstructend%
  \tagstructbegin{tag=\LBody}%
  \tagstructbegin{tag=\ProblemHeadingTag,title-o={\ProblemHeadingText}}%
  \tagmcbegin{tag=\ProblemHeadingTag,actualtext={\ProblemHeadingText}}%
  #2%
  \tagmcend%
  \tagstructend%
}
\newenvironment{examproblems}[2]{%
  \def\ProblemHeadingTag{#2}%
  \AssignTaggingSocketPlug{block/list/label}{examproblem}%
  \begin{enumerate}%
  \setcounter{enumi}{#1}%
  \renewcommand{\labelenumi}{\textbf{\theenumi.}}%
  \setlength{\topsep}{0.35em}%
  \setlength{\partopsep}{0pt}%
  \setlength{\itemsep}{0.48em}%
  \setlength{\parsep}{0.18em}%
}{\end{enumerate}}
\newenvironment{examsubparts}{%
  \AssignTaggingSocketPlug{block/list/label}{default}%
  \begin{enumerate}%
  \setlength{\topsep}{0.18em}%
  \setlength{\partopsep}{0pt}%
  \setlength{\itemsep}{0.16em}%
  \setlength{\parsep}{0.1em}%
}{\end{enumerate}}
\newenvironment{examinstructions}{%
  \par\begingroup\itshape
}{%
  \par\endgroup\vspace{0.18em}
}
\makeatletter
\renewcommand\subsection{\@startsection{subsection}{2}{0pt}%
  {-1.25ex plus -.25ex minus -.1ex}%
  {0.55ex plus .1ex}%
  {\normalfont\large\bfseries}}
\renewcommand{\@maketitle}{%
  \newpage\null\vskip -1em%
  {\footnotesize\bfseries
    \href{https://www.ufl.edu/}{University of Florida}, \href{https://math.ufl.edu/}{Department of Mathematics}\par}%
  \vskip -0.5em%
  \tagstructbegin{tag=H1,title={Algebra First-Year Practice Exam A, Part 1}}%
  \tagpdfparaOff%
  \tagmcbegin{tag=H1}%
    {\LARGE\bfseries\href{https://gma.math.ufl.edu/exams/algebra-first-year/}{Algebra First-Year Practice Exam A}, Part 1\par}%
  \tagmcend%
  \tagpdfparaOn%
  \tagstructend%
  \vskip 0.25em%
  \tagmcbegin{artifact}%
    \hrule height 1pt%
  \tagmcend%
  \vskip 1em%
}
\makeatother
\title{Algebra First-Year Practice Exam A, Part 1}
\author{}
\date{}
\begin{document}
\maketitle
\thispagestyle{empty}
\begin{examinstructions}
Answer four problems. Write your answers clearly in complete English sentences. You may quote results (within reason) as long as you state them clearly.
\end{examinstructions}
\begin{examproblems}{0}{H2}
\def\ProblemHeadingText{Problem 1.}
\item
Let~\(G\) be a group of order \(44\). Prove that
\begin{examsubparts}
\item[\textnormal{(a)}]
a Sylow \(11\)-subgroup of~\(G\) is normal in~\(G\), and
\item[\textnormal{(b)}]
the centre of~\(G\) contains an element of order~\(2\).
\end{examsubparts}
\def\ProblemHeadingText{Problem 2.}
\item
Let \(G=A\times B\) be a finite group which is the (internal) direct product of two subgroups~\(A\) and~\(B\) which are nonabelian simple groups.
\begin{examsubparts}
\item[\textnormal{(a)}]
Show that the only proper, nontrivial normal subgroups of~\(G\) are~\(A\) and~\(B\). (Hint: If~\(N\) is a proper normal subgroup different from~\(A\) and~\(B\), consider the commutator \([N,A]\).)
\item[\textnormal{(b)}]
Give an example to show that (a) would be false were~\(A\) and~\(B\) allowed to be abelian.
\end{examsubparts}
\def\ProblemHeadingText{Problem 3.}
\item
Let \(G=\mathrm{GL}_2(\mathbb{R})\) act on \(\mathbb{R}^2\) by matrix multiplication: \(A\mathbin{\cdot}\mathbf{v}=A\mathbf{v}\) for \(A\in G\), \(\mathbf{v}\in\mathbb{R}^2\).
\begin{examsubparts}
\item[\textnormal{(a)}]
Determine the stabilizer in~\(G\) of \(\mathbf{e}_2=\begin{bmatrix}0\\1\end{bmatrix}\in\mathbb{R}^2\).
\item[\textnormal{(b)}]
Determine the orbit of \(\mathbf{e}_2\).
\item[\textnormal{(c)}]
How many orbits are there for the action of~\(G\) on \(\mathbb{R}^2\)?
\end{examsubparts}
\def\ProblemHeadingText{Problem 4.}
\item
Let~\(G\) be a group.
\begin{examsubparts}
\item[\textnormal{(a)}]
Define the ascending central series
\[
Z_0(G)\leq Z_1(G)\leq Z_2(G)\leq\cdots
\]
 for~\(G\). (This is also known as the upper central series.)
\item[\textnormal{(b)}]
Prove that the groups \(Z_i(G)\) are all characteristic subgroups of~\(G\).
\item[\textnormal{(c)}]
Give an example of a nontrivial group~\(G\) such that the groups \(Z_i(G)\) are all trivial.
\end{examsubparts}
\def\ProblemHeadingText{Problem 5.}
\item
Let \(n\geq 2\) and let \(Z_n\) be a cyclic group of order~\(n\). Prove that
\[
\operatorname{Aut}(Z_n)\cong (\mathbb{Z}/n\mathbb{Z})^\times.
\]
\end{examproblems}
\end{document}
